\documentclass[11pt,letterpaper]{article}
\usepackage[T1]{fontenc}
\usepackage[utf8]{inputenc}
\usepackage{lmodern,microtype,amsmath,amssymb,amsthm,mathtools}
\ifdefined\pdfmapfile\pdfmapfile{+euler.map}\fi
\usepackage[margin=.88in,headheight=15pt,headsep=.19in,footskip=.32in]{geometry}
\usepackage{booktabs,tabularx,longtable,array,enumitem,xcolor,fancyhdr,needspace}
\usepackage{xurl}
\usepackage[colorlinks=true,linkcolor=teal!65!black,citecolor=teal!65!black,urlcolor=teal!65!black]{hyperref}
\hypersetup{pdftitle={Four copies suffice for local theta sharpness},pdfauthor={OpenAI},pdfsubject={A source-grounded simplification proposal for the Lean proof of RH and Dirichlet GRH},bookmarksnumbered=true}
\definecolor{ink}{HTML}{173246}
\definecolor{accent}{HTML}{006F73}
\definecolor{pale}{HTML}{EDF5F5}
\definecolor{quiet}{HTML}{52636D}
\pagestyle{fancy}\fancyhf{}
\fancyhead[L]{\small\sffamily\color{quiet}CM norm cancellation and local theta sharpness}
\fancyhead[R]{\small\sffamily\color{quiet}Research proposal}
\fancyfoot[R]{\small\sffamily\color{accent}\thepage}
\setlength{\parindent}{0pt}
\setlength{\parskip}{5pt plus 1pt minus .5pt}
\setlength{\emergencystretch}{2em}
\setlist{topsep=4pt,itemsep=3pt,parsep=0pt}
\renewcommand{\arraystretch}{1.12}
\raggedbottom
\clubpenalty=7000\widowpenalty=7000
\newtheorem{theorem}{Theorem}[section]
\newtheorem{lemma}[theorem]{Lemma}
\newtheorem{proposition}[theorem]{Proposition}
\theoremstyle{remark}\newtheorem{remark}[theorem]{Remark}
\newcommand{\Cp}{\mathbb C_p}
\newcommand{\Qp}{\mathbb Q_p}
\newcommand{\Zp}{\mathbb Z_p}
\newcommand{\Fp}{\mathbb F_p}
\newcommand{\OO}{\mathcal O}
\newcommand{\NN}{\mathbb N}
\newcommand{\A}{\mathcal A}
\newcommand{\Np}{\mathcal N_p}
\newcommand{\vp}{v_p}
\newcommand{\val}{\operatorname{val}}
\newcommand{\coef}[2]{[z^{#1}]\,#2}
\newcommand{\Src}[2]{\href{https://github.com/VladimirReshetnikov/ProveIt/blob/89d12dd40424f1217fd818d84877b3da52396a3b/Analysis/GeneralizedRiemannHypothesis/Lean/grh/src/#1}{#2}}
\newcommand{\PropSrc}[2]{\href{https://github.com/VladimirReshetnikov/ProveIt/blob/89d12dd40424f1217fd818d84877b3da52396a3b/Analysis/GeneralizedRiemannHypothesis/docs/proposals/#1}{#2}}
\newcommand{\decl}[1]{\textnormal{\small\texttt{\detokenize{#1}}}}
\newcommand{\keybox}[1]{\par\smallskip\noindent\colorbox{pale}{\parbox{\dimexpr\linewidth-2\fboxsep\relax}{#1}}\par\smallskip}
\newcommand{\source}[1]{\par{\small\color{quiet}\textsf{Source:} #1}\par}
\AddToHook{cmd/section/before}{\Needspace{6\baselineskip}}
\AddToHook{cmd/subsection/before}{\Needspace{4\baselineskip}}
\begin{document}
\thispagestyle{empty}
{\sffamily\bfseries\color{accent} PROOF SIMPLIFICATION\hfill 27 NOVEMBER 2026}\par
\vspace{.24in}
{\fontsize{29}{34}\selectfont\bfseries\color{ink}Four copies suffice for\\local theta sharpness\par}
\vspace{.13in}
{\Large CM norm cancellation in the Alp\"oge--Claude proof\\of RH and Dirichlet GRH\par}
\vspace{.18in}
{\large Analysis of the supplied Lean development}\par
\vspace{.18in}
\keybox{\textbf{Recommendation.} Replace cancellation by squares of rational
$p$-adic scalars with cancellation by local CM norms. This reduces the local
sharpness construction from twelve copies to four, with \emph{the same valuation
bound}. First recover the existing twelve-copy interface by zero-degree padding;
then consider a separate geometric refactor.}

\textbf{Main result of the analysis.} Under the same hypotheses as the
source's twelve-copy theorem, the paper argument below gives
\begin{equation*}
 \boxed{\vp(l!)+\val_4(l)\ \leq\ S_p+2\sum_\tau\vp(B_\tau).}
\end{equation*}
Here $S_p$ and $B_\tau$ are the \emph{unchanged} local prime sum and loss
products of the existing proof. The notation $\val_n(l)\leq r$ is an
existential coefficient bound, defined precisely in Section~\ref{sec:notation}.

\begin{center}
\begin{tabular}{lll}
\toprule
Prime case & New construction & Copies needed\\
\midrule
Split, including $2$ & Norm pair $(1,-1)$ & $2$\\
Inert, $p>k$ & Two coupling copies; one spanning copy & $3$\\
Inert, $p\leq k$ & Two coupling copies; two residual copies & $4$\\
Ramified, $p=d$ & Three rational scales; no near pairs & $3$\\
\bottomrule
\end{tabular}
\end{center}

\textbf{Status.} This is a mathematical replacement argument from named
theorems in the pinned source, with independently inspected source interfaces,
an exclusion audit of the existing proposals, and exact finite sanity checks.
The supplied Lean module is explicitly \textbf{uncompiled}; it is a candidate
for part of the refactor, not a completed formalization of the four-copy theorem.
The original development was not rebuilt or kernel-replayed in this session.

\textbf{Source revision.} All repository assertions refer to ProveIt commit
\begin{center}\small\ttfamily 89d12dd40424f1217fd818d84877b3da52396a3b\end{center}
The accompanying archive contains this \texttt{.tex}, its compiled PDF,
candidate Lean code, exact-check scripts and results, and immutable source
inventories with Git blob hashes.
\newpage
\begingroup\small\setlength{\parskip}{0pt}\tableofcontents\endgroup
\newpage

\section{What was analyzed and what is being proposed}
\label{sec:scope}

The supplied development proves Mathlib's \decl{RiemannHypothesis} and
nonvanishing to the right of $\operatorname{Re}s=1/2$ for Dirichlet
$L$-functions, with the usual exclusion of the principal-character pole.
Its public zero-on-the-line statements cover every positive modulus and
primitive or imprimitive Dirichlet characters. In this report, ``GRH'' means
\emph{this Dirichlet GRH}; no claim about every automorphic or Artin
$L$-function is being inferred. The public statement file imports only
Mathlib, and the public proof file connects those statements to
\decl{New.grh}.\cite{public}

The selected Lean source tree contains 803 files and 477,124 lines, including
comments and blank lines. It pins Lean 4.35.0 and Mathlib commit
\texttt{56a005855ea17e0b0958f73ec801ae140b0a1679}. Its source-tree hash is
\begin{center}\small\ttfamily 55fe9c6cb2e5be421c5871e772b6d736ecb5d164.\end{center}
That tree agrees with the readable edition of upstream
\texttt{lsernau/grh-lean} at commit
\href{https://github.com/lsernau/grh-lean/commit/72fead102f663d5c7a407964a13dc0cf9d977207}{\texttt{72fead102f}}.
Andrew Sutherland's verification page reports compilation, replay of all
803 modules, independent checker acceptance, and a successful paranoid
comparator run for that edition. These are \emph{external verification
results}, not a replay performed for this report.\cite{verification}

The review followed the reduction to the main determinant inequality, the
local sharpness-to-global-form bridge, the local formal-group and CM
integrality arguments on which that bridge relies, and all existing proposal
manuscripts and plans. It does not claim an independent line-by-line audit
of the entire development. The detailed contribution is confined to a
specific arithmetic step whose interfaces can be inspected precisely.

\subsection{Why choose this simplification}

Many natural analytic and linear-algebra simplifications are already
proposed in the supplied directory. The local derivative sharpness argument,
in contrast, retains a uniform six-copy construction based on four squares.
An existing CM integrality theorem is stronger than the rational-square
specialization used there. Applying its full strength supplies a concrete
replacement, rather than another suggestion to optimize an asymptotic
constant or to repackage the endpoint.

The first implementation should replace the proof of a local theorem while
retaining its external type. This makes the simplification compatible with
the existing proposals: improvements to spanning, auxiliary-prime choice,
torsion interpolation and cutoff arguments can be implemented independently.
No saving in total source size or build time is asserted before an actual
formalization and measurement.

\section{Where local sharpness enters the proof}
\label{sec:architecture}

The central comparison bounds one corrected covolume from above and below.
A short integral source family and sufficiently many Poincar\'e translates
give an upper bound. Arithmetic height and denominator control give a lower
bound. With $D$ denoting admissible data, the formal conclusion is
\begin{equation}
 D.m\,D.\mathrm{totalRemainder}
 \leq D.f\,D.a\,(D.\mathrm{logB})^5.
 \label{eq:main}
\end{equation}
The same large main term occurs on both sides and cancels. The finite index
and the cost of completing the selected family are included in the
comparison; they cannot simply be dropped.\cite{mainineq}

The admissible lift-chain length satisfies $\sqrt{x}\leq m\leq2\sqrt{x}$.
The exact relation between the total remainder and the two prime-sum cutoffs,
together with a prime number theorem at the distant cutoff, converts
\eqref{eq:main} to
\begin{equation}
 R_{q,d}(x)\leq x^{-1/2+\varepsilon}
 \quad\hbox{for all sufficiently large }x.
 \label{eq:onesided}
\end{equation}
Here $R_{q,d}$ is the centered smoothed sum with nonnegative weights
$\Lambda(n)(1-\chi_{-d}(n))$ on $n\equiv1\pmod q$.
The auxiliary prime $d$ is fixed before $x$ tends to infinity. A suitable
choice of $d$ removes twisted zeros from the relevant rectangle, so that
Landau's sign-sensitive one-sided argument applies; functional-equation
reflection then gives the critical line.\cite{endpoint}

The arithmetic lower bound uses a chain-to-single-lift factorization,
comparison operators, special values and denominator estimates from
congruences and Selmer groups. The local sharpness calculation supplies
coefficient witnesses for rank-one theta test vectors. It is the existence
of a witness at a particular valuation, not merely a divisibility estimate
in the opposite direction, that is needed.

This distinction explains the target of the proposal. Weakening the local
leading coefficient would change the main arithmetic term and can destroy
the cancellation in \eqref{eq:main}. Reducing the number of copies
\emph{without changing the valuation budget} is safe at that interface.
The new argument keeps the normal form, the coupling coefficient identity,
the factorial losses, and the eventual main term intact.\cite{couplings,cases,boundattained}

\subsection{An analytic shortcut that the review does not endorse}

For a fixed quadratic twist, the relevant transform detects the difference
of untwisted and twisted zero multiplicities. A zero can therefore cancel,
and a twisted real zero can introduce a real residue with the wrong sign.
An upper bound alone cannot rule out this possibility: the model remainder
$-x^{\alpha-1}$ has the desired upper bound for smaller exponents but has a
negative-residue pole at $\alpha$. Thus replacing the source's twist selection
with an arbitrary fixed twist is not a valid alternative proof. The existing
guide and proposal corpus already discuss this obstruction.\cite{endpoint,guide,sparse}

\section{Exclusion audit: how this differs from the existing proposals}
\label{sec:exclusions}

The audit covered all five bundles and twelve TeX sources, including
included sections, README files, integration plans, and explicitly proposed
future work. The inventories in \texttt{sources/} identify every audited
file by its pinned path and blob hash. The exclusion criterion was broader
than ``already implemented'': an idea explicitly proposed for future work
was treated as occupied.

\begin{longtable}{@{}p{.25\linewidth}p{.69\linewidth}@{}}
\toprule
Bundle & Ideas treated as already proposed\\
\midrule
\endhead
\PropSrc{centered_theta_lift_simplification/centered_theta_lifts.tex}{Centered theta lifts}
& Centered covariance, adaptive rank completion, random contraction products,
continuous rank--height budgets, weaker moment tolerances, charging index by
new rank, and partial/exterior-power determinant arguments.\\[4pt]
\PropSrc{proof_simplifications/proof_simplifications.tex}{Proof simplifications}
& Direct ramified-source rank recurrence; short Dirichlet polynomials and
mollifiers; quadratic large sieve or Fourier completion; direct smoothed
Perron estimates at the distant cutoff.\\[4pt]
\PropSrc{rh-simplifications/rh-simplifications.tex}{RH simplifications}
& Rank recurrence and its obstruction; simultaneous torsion interpolation,
including powers of a common character; large-sieve twist choice; endpoint
assembly and further interpolation-grid refinements.\\[4pt]
\PropSrc{rh_sparse_theta_proof/sparse_theta_proof.tex}{Sparse theta proof}
& Reciprocal convexity, sparse fourth-power and compatible dyadic grids,
window averages, Gaussian theta-tail second moments, divergent auxiliary
prime subfamilies, data reuse across blocks, and averaged volume bounds.\\[4pt]
\PropSrc{shorter-theta-lift-proof/shorter_theta_proof.tex}{Shorter theta-lift proof}
& Centered spanning, elementary quadratic mean squares and Fourier
completion, a large-sieve alternative, fixed-modulus prime number theorem
shortcuts, and monotone ceiling interpolation.\\
\bottomrule
\end{longtable}

None of these manuscripts or plans proposes using
\decl{normPairClosure} to replace the six rational-square scales, the
inert norm-minus-one root of unity below, the split limit pair $(1,-1)$,
or the resulting four-copy local sharpness theorem. Their occurrences of
``sharpness'' concern other estimates or the retention of the derivative
input. The torsion-interpolation proposal chooses specializations for
Selmer/special-value data; it does not cancel norm weights in a local theta
product.

The novelty claim is consequently precise and limited: this is a distinct
proposal \emph{relative to the user's specified corpus at the pinned
revision}. No claim of priority over all mathematical literature is made.

\section{Notation and the exact coefficient interface}
\label{sec:notation}

Fix a prime $d\equiv7\pmod8$, the CM field
$K=\mathbb Q(\sqrt{-d})$, a prime $p$, and the source's chosen embedding
$\iota_p$ into $\Cp$. Set $\OO=\OO_{\Cp}$. We normalize
$\vp(p)=1$, so $|a|_p=p^{-\vp(a)}$ for nonzero $a$; negative valuations
are allowed. For a multi-index $\nu=(\nu_x)$, put
\[
 |\nu|=\sum_x\nu_x,\qquad
 \nu!=\prod_x\nu_x!,\qquad
 V_p(\nu)=\vp(\nu!)=\sum_x\vp(\nu_x!).
\]

Let $C$ be the source CM lattice, let $\theta(0)=1$ be its normalized
translated theta series, and write
\[
 g(z)=-\log\theta(z),\qquad g(0)=0,\qquad\theta(z)=\exp(-g(z)).
\]
All statements here concern formal series. The source's parameter
$y=\mathrm{translatedY}$ has $y(0)=0$ and linear coefficient $1$. Define
\[
 \A=\OO[[y(z)]]\subset\Cp[[z]].
\]
This denotes the image of integral power series under substitution by
$y(z)$, namely \decl{C.integralSeriesRing}; it need not be identified
with the ring of series integral in the original coordinate $z$.

For the finite family of linear forms $\mu_j(z)$ and nonnegative
multiplicities $c_j$, set
\[
 \Theta(z)=\prod_j\theta(\mu_j(z))^{c_j},\qquad
 \A_\mu=\OO[[\,y(\mu_j(z)):j\in J\,]],
\]
where $\A_\mu$ is the source's substituted multivariate series ring.
In particular $\Theta(0)=1$, whether or not $\Theta\in\A_\mu$.

The assertion $\val_n(\nu)\leq r$ will mean exactly that there exist
$H_i\in\A_\mu$ and multi-indices $\alpha_i$ for $0\leq i<n$ such that
\begin{equation}
 \sum_{i<n}\alpha_i=\nu,\qquad
 p^{-r}\leq\prod_{i<n}\bigl|[z^{\alpha_i}](H_i\Theta)\bigr|_p.
 \label{eq:leastval}
\end{equation}
This is \decl{C.LeastValLE mu c n nu r}. The notation is shorthand for an
existential proposition; it does not assert that an infimum is attained
or introduce a new real-valued function.

Two source rules will be used repeatedly. Bounds add under concatenating
independent copies, and extra degree-zero copies have zero valuation cost:
\begin{align}
 \val_n(\nu)\leq r,\ \val_{n'}(\eta)\leq s
 &\ \Longrightarrow\ \val_{n+n'}(\nu+\eta)\leq r+s,\label{eq:add}\\
 \val_n(\nu)\leq r
 &\ \Longrightarrow\ \val_{n+h}(\nu)\leq r\quad(h\geq0).
 \label{eq:pad}
\end{align}
For the second rule take the new $H_i=1$ and $\alpha_i=0$.
\source{\Src{New/Derivatives/Sharpness/LeastValuation.lean}{LeastValuation},
lines 259--306: \decl{LeastValLE}, \decl{leastValLE_zero}, and
\decl{LeastValLE.add}.}

\section{The unused generality: cancellation by CM norms}
\label{sec:normbalance}

Define the closed set of local norm pairs
\begin{equation}
 \Np=\overline{\{(\iota_p(\alpha),\iota_p(\alpha\bar\alpha)):
                         \alpha\in\OO_K\}}\ \subset\Cp^2.
 \label{eq:normclosure}
\end{equation}
This is exactly \decl{normPairClosure d}. The source proves
\begin{equation}
  (\delta,s)\in\Np
  \quad\Longrightarrow\quad
  \exp\!\bigl(sg(z)-g(\delta z)\bigr)\in\A.
  \label{eq:cm_input}
\end{equation}
For an actual CM integer this follows from its theta isogeny relation.
The theorem for limits follows by coefficientwise continuity and closedness
of the integral-series ring. These are existing, proved source inputs.
\source{\Src{Lib/General/NumberTheory/PadicLimitsValuation.lean}{PadicLimitsValuation},
line 169; \Src{Lib/CM/Integrality/LogTheta.lean}{LogTheta}, lines 271--318,
especially \decl{expSubst_sub_rescale_mem}.}

\begin{lemma}[Norm-balanced theta product]\label{lem:balance}
Suppose $(\delta_i,s_i)\in\Np$ for $0\leq i<n$ and
$\sum_i s_i=0$. Then
\[
 u(z)=\prod_{i<n}\theta(\delta_i z)\in\A^\times,\qquad u(0)=1.
\]
Consequently $U(z)=\prod_{i<n}\Theta(\delta_i z)$ is a unit in $\A_\mu$
with constant term $1$.
\end{lemma}

\begin{proof}
All exponential arguments have zero constant term. Hence formal
exponentials multiply according to addition of their arguments, giving
\begin{align*}
 \prod_i\exp(s_i g(z)-g(\delta_i z))
 &=\exp\!\left(\Bigl(\sum_i s_i\Bigr)g(z)-\sum_i g(\delta_i z)\right)\\
 &=\prod_i\theta(\delta_i z).
\end{align*}
Each factor on the left belongs to $\A$ by \eqref{eq:cm_input}.
The product has constant term $1$. Writing it as a series in $y$, its
formal inverse is again an integral series in $y$ (the coefficient
recursion divides only by the unit constant term). Let $v\in\A$ satisfy
$uv=1$.

Substitution in $\mu_j(z)$ gives
\[
 U(z)=\prod_j u(\mu_j(z))^{c_j},\qquad
 U(z)^{-1}=\prod_j v(\mu_j(z))^{c_j}\in\A_\mu.
\]
Thus inversion is justified in one variable before substitution. No
general assertion about units of an arbitrary subring is needed.
\end{proof}

In particular $(1,1)\in\Np$ and any $(\delta,-1)\in\Np$ give
\begin{equation}
 \theta(z)\theta(\delta z)\in\A^\times,
 \qquad \Theta(z)\Theta(\delta z)\in\A_\mu^\times.
 \label{eq:two}
\end{equation}
The current proof restricts to $(a,a^2)$ with $a\in\Zp$. It represents
$-1$ by four squares, uses the scales $(1,a_1,a_2,a_3,a_4,0)$, and obtains
six factors. The last zero is padding. Lemma~\ref{lem:balance} replaces
this restriction by the natural CM norm coordinate.\cite{leastval}

\section{Local constructions with two or three factors}
\label{sec:local}

\subsection{Inert odd primes: a root of unity of norm minus one}

Assume $p$ is inert in $K$ and odd. Choose $\zeta\in\Cp$ of exact order
$2(p+1)$. Since
\[
 2(p+1)\mid p^2-1,
\]
we have $\zeta^{p^2-1}=1$ and $|\zeta|_p=1$. Also
$\zeta^{p+1}$ has order two, so it equals $-1$ in characteristic zero.
The source theorem
\decl{mem_normPairClosure_of_pow_eq_one} states, in this inert case,
\[
 \zeta^{p^2-1}=1
 \quad\Longrightarrow\quad
 (\zeta,\zeta^{p+1})\in\Np.
\]
Thus $(\zeta,-1)\in\Np$, and \eqref{eq:two} gives the required
two-factor unit.

The multipliers are $1$ and $\zeta$, both of norm one. There is no need
for $\zeta\in\Qp$ or $\zeta^2=-1$. The scalar norm used by the CM
identity is $\zeta^{p+1}$, not $\zeta^2$.
Since $d\equiv7\pmod8$, $2$ splits in $K$; moreover the admissible data
in the application force every nonsplit prime to be at least $13$.
There is consequently no inert-$2$ case in the theorem being replaced.
\source{\Src{Lib/CM/FormalGroup/FormalParameter.lean}{FormalParameter}, line 44;
\Src{New/Reduction/Data/FieldsAndPlaces.lean}{FieldsAndPlaces}, line 162.}

\subsection{Split primes, including two: the closure point \texorpdfstring{$(1,-1)$}{(1,-1)}}

Let $p\OO_K=\mathfrak p\bar{\mathfrak p}$ with distinct prime ideals,
where $\mathfrak p$ is the place selected by $\iota_p$. For every
$n\geq1$, the Chinese remainder theorem supplies $\alpha_n\in\OO_K$
with
\[
 \alpha_n\equiv1\pmod{\mathfrak p^n},\qquad
 \alpha_n\equiv-1\pmod{\bar{\mathfrak p}^{\,n}}.
\]
In the chosen embedding, $\iota_p(\alpha_n)\longrightarrow1$ and
$\iota_p(\bar\alpha_n)\longrightarrow-1$. Multiplication is continuous
in $\Cp$, so their norm pairs converge to $(1,-1)$. Hence
\begin{equation}
 (1,-1)\in\Np,\qquad\theta(z)^2\in\A^\times,
 \qquad\Theta(z)^2\in\A_\mu^\times.
 \label{eq:split}
\end{equation}

This argument applies at $2$: the two primes above $2$ are still distinct.
It does not divide by $2$ in the coefficient ring. It also does not
assert the existence of a global CM integer or unit of norm $-1$;
the point is obtained only in a $p$-adic closure.

A second, elementary description is useful for implementation and checking.
Let $b=\iota_p(\sqrt{-d})\in\Zp^\times$, which lies in $\Qp$ at a
split prime. Choose ordinary integers $n_j\to b^{-1}$ and take
$\alpha_j=n_j\sqrt{-d}$. Then
\[
 \iota_p(\alpha_j)\to1,\qquad
 \iota_p(\alpha_j\bar\alpha_j)=d n_j^2\to d/b^2=-1.
\]
At $2$, the needed square root exists because $-d\equiv1\pmod8$.
The CRT proof avoids introducing that square-root lemma separately.

This local-density result must be added to Lean: it is not a consequence
of the existing Teichm\"uller-pair theorem alone. At $p=2$, membership
of $\theta^2$ does not justify membership of $\theta$ by dividing an
integral square-root recursion by $2$; two copies are retained here.
\source{\Src{Lib/CM/FormalGroup/PrimeNotation.lean}{PrimeNotation}, line 352,
\decl{QuadField.star_ne_of_chiK_eq_one}, supplies the distinct-place fact.}

\subsection{The ramified prime: three rational scales}

Here $p=d\equiv3\pmod4$. In $\Fp$, the two sets
\[
 \{a^2:a\in\Fp\},\qquad\{-1-b^2:b\in\Fp\}
\]
each have $(p+1)/2$ elements, so they intersect. Thus
$a_0^2+b_0^2=-1$ has a solution. Both coordinates are nonzero, since
$-1$ is not a square in $\Fp$. Fix an integer lift $b$ of $b_0$.
Hensel's lemma, applied to $X^2+b^2+1$ at $a_0$, supplies
$a\in\Zp$ with
\begin{equation}
 1+a^2+b^2=0.
 \label{eq:ramified}
\end{equation}
The derivative $2a_0$ is a unit, so this is an ordinary simple-root lift.

Density of integers in $\Zp$ gives $(t,t^2)\in\Np$ for every
$t\in\Zp$. Applying Lemma~\ref{lem:balance} to
$(1,1),(a,a^2),(b,b^2)$ proves
\[
 \theta(z)\theta(az)\theta(bz)\in\A^\times
\]
and the corresponding three-factor statement for $\Theta$.

There is a precise obstruction to using the same two-factor recipe here.
Modulo the ramified prime, conjugation acts trivially on the residue field;
the residue of a norm is therefore a square. It cannot be $-1$.
Consequently $(\delta,-1)\notin\Np$ for all $\delta$ at this prime.
Three is minimal for this \emph{norm-balancing recipe with the first pair
$(1,1)$}. This is not a lower bound for every conceivable sharpness proof.

Crucially, near pairs in the source are necessarily inert. At $p=d$ there
are no near pairs, hence no second coupling construction to combine with
these three factors.\cite{nearpairs}

\section{Transferring cancellation to a coefficient witness}
\label{sec:transfer}

\begin{lemma}[Copy-independent transfer]\label{lem:transfer}
Suppose $\delta_0=1$, $|\delta_i|_p\leq1$, and
$U=\prod_{i<n}\Theta(\delta_i z)\in\A_\mu^\times$.
For every $G\in\A_\mu$, a bound
$|[z^\nu]G|_p\geq p^{-r}$ implies $\val_n(\nu)\leq r$.
\end{lemma}
\begin{proof}
Take $H_0=GU^{-1}$ and $H_i=1$ for $i>0$. Because $\delta_0=1$,
\[
 \prod_{i<n}(H_i\Theta)(\delta_i z)=G(z).
\]
Its coefficient at $\nu$ is a finite sum over $\sum_i\alpha_i=\nu$ of
\[
 \left(\prod_i\delta_i^{|\alpha_i|}\right)
 \prod_i[z^{\alpha_i}](H_i\Theta).
\]
The first factor has norm at most one. The ultrametric inequality gives
one summand with norm at least $p^{-r}$, and therefore the required product
in \eqref{eq:leastval} has norm at least $p^{-r}$.
\end{proof}

This is the existing \decl{leastValLE_of_rescale}; that theorem already
allows arbitrary integral scalars in $\Cp$. Its type does not restrict
the scales to $\Zp$. Only the current constructor for a cancelling product
has that more restrictive type.\cite{leastval}

\subsection{The same Katz bound with fewer copies}

Define the source scaling exponent and derivative rate by
\[
\begin{array}{c|ccc}
&\text{split}&\text{inert}&\text{ramified}\\ \hline
\rho_p&1/(p-1)&1/(p^2-1)&1/(2p-2)\\[2pt]
\gamma_p=1/(p-1)-\rho_p&0&p/(p^2-1)&1/(2p-2).
\end{array}
\]
Assume some linear form $\mu_{j_0}$ has a unit coefficient in every
coordinate. The source Katz divisibility-limit theorem supplies
$h\in\A$ with
\[
 \vp\bigl(h^{(|\nu|)}(0)\bigr)\leq\gamma_p|\nu|.
\]
For $G(z)=h(\mu_{j_0}(z))$, the multinomial formula gives
\[
 [z^\nu]G=\frac{h^{(|\nu|)}(0)}{\nu!}
                   \prod_x\mu_{j_0,x}^{\nu_x}.
\]
All coefficients in the last product are units. Thus
\[
 \vp([z^\nu]G)\leq\gamma_p|\nu|-V_p(\nu).
\]
Apply Lemma~\ref{lem:transfer} with the products constructed above:
\begin{equation}
 \val_{c(p)}(\nu)\leq\gamma_p|\nu|-V_p(\nu),\qquad
 c(p)=\begin{cases}2,&p\ne d,\\3,&p=d.\end{cases}
 \label{eq:katz}
\end{equation}
This is the proof of \decl{leastValLE_of_unit_coefficients} with a
different cancelling product. No new Katz estimate is needed.

\subsection{The one-copy spanning witness is already in the source}

Suppose each coordinate can be recovered from the forms with integral
coefficients and a common unit $s$:
\[
 \sum_j a_{xj}\mu_j(z)=s z_x,\qquad a_{xj}\in\OO,\quad |s|_p=1.
\]
Then $G_x=\sum_j a_{xj}y(\mu_j(z))\in\A_\mu$ has leading homogeneous
term $s z_x$. Since $\Theta(0)=1$, the coefficient at $z^\nu$ in
\[
 \left(\prod_xG_x^{\nu_x}\right)\Theta
\]
is $s^{|\nu|}$, a unit. Consequently
\begin{equation}
 \val_1(\nu)\leq0.
 \label{eq:span}
\end{equation}
The existing proof of \decl{leastValLE_zero_of_span} constructs exactly
this assertion as \decl{h_one_copy}, then adds five zero-degree copies.
Exposing the stronger statement is an interface change, not a new
mathematical ingredient.\cite{leastval}

\section{The coupling estimate also needs only two copies}
\label{sec:coupling}

At an inert prime, the existing formal-group analysis writes the scaled
theta product as
\[
 F(z)=\Theta(\varpi z)=W(z)\exp(\kappa\Xi(z)),\qquad
 \vp(\varpi)=\rho_p,\qquad \kappa=\varpi^{p+1}/p,
\]
with $W$ an integral unit and
$\vp(\kappa)=1/(p-1)-1<0$. The series $\Xi$ and its unit coefficients
satisfy the existing normal-form hypotheses. A coupling family supplies
degrees $\nu',\nu''$ and multiplicities $j_\tau$, extended by zero outside
its set $\Phi_0$ of coupling starts, with
\[
 \nu=\nu'+\nu'',\qquad J=\sum_\tau j_\tau,\qquad
 |\nu''|=J,\qquad |\nu'|\geq pJ.
\]

The decisive source theorem,
\decl{IsCouplingFamily.norm_coeff_pairExponent}, is already stated for
an arbitrary number $n+1$ of copies. It requires $\delta_0=1$ and an
integral product $\prod_iF(\delta_i z)$ of constant term $1$.
It does not require six copies or rational $p$-adic scales. The two-factor
CM product supplies this hypothesis after scaling, because the source
maps $\A_\mu$ to integral series under $z\mapsto\varpi z$.
\cite{couplings}

For $\delta=(1,\zeta)$ the theorem applies to
\[
 F(z'+z'')F(\zeta z').
\]
It states that the valuation of its coefficient at
$(z')^{\nu'}(z'')^{\nu''}$ is
$J\vp(\kappa)-V_p(\nu'')$. Undoing the common $\varpi$ scaling and
applying the integral-substitution coefficient extraction gives
\begin{equation}
 \val_2(\nu)\leq
 \frac{J}{p-1}-J-\rho_p|\nu|-V_p(\nu'').
 \label{eq:couplingprecise}
\end{equation}
From $|\nu|\geq(p+1)J$ and $\rho_p(p+1)=1/(p-1)$, the first
useful consequence is
\begin{equation}
 \val_2(\nu)\leq-J.
 \label{eq:couplingzero}
\end{equation}
Legendre's inequality $V_p(\nu')\leq|\nu'|/(p-1)$ gives the other:
\begin{equation}
 V_p(\nu')+V_p(\nu'')+\val_2(\nu)
 \leq\gamma_p|\nu|-J.
 \label{eq:couplinggamma}
\end{equation}
These are exactly the numerical bounds in the current six-copy wrapper.
Only the copy count has changed.

\begin{remark}[A check directly in the normal form]
The monomial degrees in the inert coupling series are
$(p^a+1)p^b$ with $a$ odd. Modulo $2(p+1)$ this is $p+1$:
for odd $p,a$, the quotient $(p^a+1)/(p+1)$ is odd, and $p^b$ is odd.
Thus $\zeta^{(p^a+1)p^b}=-1$ and $\Xi(\zeta z)=-\Xi(z)$.
It follows directly that $F(z)F(\zeta z)=W(z)W(\zeta z)$ is integral.
This is a useful consistency check. The CM-norm argument additionally
places the unscaled product in $\A_\mu$, which is what the residual
Katz argument requires.
\end{remark}

\section{The four-copy replacement theorem}
\label{sec:four}

Keep the exact admissible data $D$, CM lattice, resolution forms,
multiplicities, and near-weight condition of
\decl{Data.leastValLE_twelve}. Explicitly, the frame is tight, one
resolution form has unit coefficients at the chosen place, and the
multi-index satisfies
\[
 \kappa_\tau\leq l_\tau\leq\kappa_\tau+x^{42}.
\]
Let
\begin{align*}
 n^\sharp&=\lceil6b/k\rceil, & l_\tau^-&=\max(l_\tau-n^\sharp,0),\\
 B_\tau&=\prod_{j=1}^{n^\sharp}(l_\tau^-+j), &
 L_p(l)&=\sum_\tau\vp(B_\tau),\\
 S_p&=D.\mathrm{localPrimeSumTotal}(p), &
 E_p(l)&=S_p+2L_p(l).
\end{align*}
The finite sums over $\tau$ suppress the source's harmless
$\Phi\times\mathrm{Unit}$ indexing. The $B_\tau$ are positive integers,
so $L_p(l)\geq0$. These definitions are unchanged.\cite{nearpairs,cases}

More explicitly, the local prime sum is
\[
 S_p=a_p\sum_\tau\sum_{e\in I_p}
       \left(\frac{\kappa_\tau}{p^e}-\kappa_{u(-p^e)\tau}\right)_+,
 \qquad
 a_p=\begin{cases}1/2,&p=d,\\1,&p\ne d.\end{cases}
\]
Here $u(-p^e)\tau$ denotes the source's partner action, and $I_p$ is empty
at split primes, consists of positive odd integers at inert primes, and
of all positive integers at the ramified prime. There is no $\log p$
factor in this local sum.
\source{\Src{New/UpperBound/Size/ShortColumns.lean}{ShortColumns}, lines 72--78.}

\begin{theorem}[Four-copy local sharpness]\label{thm:four}
Under exactly the hypotheses of \decl{Data.leastValLE_twelve},
\begin{equation}
 \val_4(l)\leq E_p(l)-V_p(l).
 \label{eq:four}
\end{equation}
The theorem is a paper-level consequence of the stated existing inputs
and the local constructions in this report. Its new Lean implementation
is an outstanding formalization task.
\end{theorem}

\begin{proof}
Use the same prime cases as the source, separating the ramified prime from
the small nonsplit case.

\textbf{Split primes.}
Here $\gamma_p=0$ and $S_p=0$.
Equation~\eqref{eq:katz} yields $\val_2(l)\leq-V_p(l)$.
Add two zero-degree copies and use $L_p(l)\geq0$ to obtain
\eqref{eq:four}. This includes $p=2$.

\textbf{Inert primes $p>k$.}
The source's coupling family and truncation decompose $l=\nu+\eta$.
Its unchanged factorial estimate gives
\[
 V_p(l)\leq J+S_p+2L_p(l).
\]
The coupling witness \eqref{eq:couplingzero} costs two copies and
valuation $-J$. The tight-frame spanning argument valid in this branch
supplies \eqref{eq:span} for $\eta$, costing one copy and valuation zero.
Their concatenation gives
\[
 V_p(l)+\val_3(l)\leq V_p(l)-J\leq S_p+2L_p(l).
\]
Pad by one copy. The assertion that a nonsplit prime $p>k$ is inert
uses the original admissible inequality $d\leq k$.

\textbf{Inert primes $p\leq k$.}
Use the original set of near-pair starts and its exponents $e_\tau$.
Set $j_\tau=\lfloor l_\tau/p^{e_\tau}\rfloor$ at those starts and
$j_\tau=0$ elsewhere.
The source decomposes
$l=\nu'+\nu''+\eta$, with $\nu'+\nu''$ the coupling degree.
Combine the two-copy bound \eqref{eq:couplinggamma} with the
two-copy residual bound \eqref{eq:katz}. By \eqref{eq:add},
\begin{equation}
 \val_4(l)\leq\gamma_p|l|-J
                     -V_p(\nu')-V_p(\nu'')-V_p(\eta).
 \label{eq:smalladd}
\end{equation}
Put $C_0=3f(\log b+2)$.
Two existing numerical estimates, with no copy count in their statements,
are
\begin{align}
 V_p(l)&\leq V_p(\nu')+V_p(\nu'')+V_p(\eta)+C_0,
 \label{eq:factorialloss}\\
 \gamma_p|l|-J+C_0&\leq S_p+L_p(l).
 \label{eq:nearcompare}
\end{align}
Combining them with \eqref{eq:smalladd} gives the slightly stronger
$V_p(l)+\val_4(l)\leq S_p+L_p(l)$, which implies the stated bound.

\textbf{The ramified prime $p=d$.}
Here $d\leq k$. A near pair would force inertness, so the near-pair
set is empty. Take $j=0$ and any unused exponent assignment. The same
comparison estimate \eqref{eq:nearcompare} applies vacuously to this
empty family and becomes
\[
 \gamma_p|l|+C_0\leq S_p+L_p(l).
\]
The three-copy bound \eqref{eq:katz} applies directly to all of $l$;
there is no coupling block. Since $C_0\geq0$,
\[
 V_p(l)+\val_3(l)\leq\gamma_p|l|\leq S_p+L_p(l).
\]
Pad by one copy and use $L_p(l)\geq0$.
\end{proof}

\source{\Src{New/Derivatives/Sharpness/Cases.lean}{Cases}, lines 403,
526--579, 619--690; \Src{New/Derivatives/Sharpness/NearPairs.lean}{NearPairs},
\decl{IsNearPair.ne_d} and \decl{IsNearPair.inert}.}

\subsection{Why the same argument does not establish three copies uniformly}

In the small inert case the residual valuation bound must be
$\gamma_p|\eta|-V_p(\eta)$, which can be negative. For one variable
with $\eta=p^2$, Legendre's formula gives $V_p(\eta)=p+1$, while
\[
 \gamma_p p^2-(p+1)=\frac{p}{p^2-1}-1<0.
\]
The one-copy spanning lemma yields only a bound of zero, so it cannot
replace this residual estimate. The present proof establishes a uniform
four-copy result; it neither establishes nor rules out a different
uniform three-copy argument.

\section{Why the local multipliers do not create global height costs}
\label{sec:global}

The scalars $\zeta$, $a$, and $b$ may depend on $p$ and the chosen
embedding. This is compatible with the actual quantifiers. They are
temporary devices for proving \eqref{eq:leastval}, not coefficients that
must be placed in one global theta test function.

The source theorem \decl{LeastValLE.exists_monomial} replaces each
arbitrary $H_i\in\A_\mu$ by a monomial
\[
 \prod_j y(\mu_j(z))^{e_{ij}}
\]
without worsening the relevant coefficient norm. Only finitely many
monomials can contribute below the specified degree. Expanding each
integral series and taking an ultrametric maximum selects one; its
integral coefficient has norm at most one and can be removed.
The exponents are natural numbers and can be bounded in terms of $|l|$.
The locally selected CM scales have disappeared by this stage.
\cite{leastval}

The following \decl{BoundAttained} argument realizes those fixed-function
monomials as algebraic theta vectors, then uses finite Fourier inversion
to select an indicator datum. The datum may depend on $p$ and the
embedding, as it already does in the original proof.
The common global scalar is chosen \emph{before} the universal quantifier
over primes and embeddings. In the existing twelve-copy implementation
it has the shape
\[
 t=M^{12f}c_0\prod_\tau B_\tau^2.
\]
Proving the four-copy bound and padding it to twelve preserves that
choice and the conjugate-height bound literally. In particular, the
proposal adds no global field containing all the local $\zeta$ and asks
for no simultaneous global norm-minus-one element.\cite{boundattained}

This is an essential compatibility check. If the local scales had survived
as coefficients in the final global test form, their arithmetic and
archimedean costs would need a new analysis. The existing extraction
interface is precisely what makes that extra task unnecessary here.

\section{A concrete Lean integration plan}
\label{sec:implementation}

\subsection{Stage one: replace the local proof and preserve its public interface}

The recommended first change consists of the following finite sequence.
\begin{enumerate}
\item Add the general norm-balanced product lemma next to the current
square-balanced lemma. Its proof follows the same formal exponential
identity and calls \decl{expSubst_sub_rescale_mem} for each pair.
\item Add selectors for an inert norm-minus-one root of unity, the split
closure point $(1,-1)$, and the ramified three-scale construction.
The split-density and two-squares/Hensel arguments need new local lemmas.
\item Generalize the scaled theta product and its unit witnesses to
integral $\Cp$ scalars and a variable number of copies. Generalize the
unit-coefficient wrapper over such a witness; do not duplicate the Katz
proof for each prime case.
\item Generalize the coupling wrapper over the number of copies. Its
central coefficient theorem is already generic. Instantiate with two
copies in the inert case.
\item Expose the one-copy result already present in
\decl{leastValLE_zero_of_span}; retain the old six-copy statement as a
padded corollary if useful to avoid intermediate breakage.
\item Prove \decl{Data.leastValLE_four}, using the four cases in
Theorem~\ref{thm:four} and the original numerical estimates.
\item Derive \decl{Data.leastValLE_twelve} by adding eight degree-zero
copies with \decl{LeastValLE.add} and \decl{leastValLE_zero}.
\end{enumerate}

The result of this stage is a simplification of the local proof and a
stronger intermediate theorem. \decl{BoundAttained}, the common
normalizing scalar, and the downstream determinant and analytic theorems
can retain their existing statements and proof bodies.
The helper \decl{padicInt_sum_four_squares} is used only by the
six-scale constructor within the selected source tree, so that particular
four-square/compactness route becomes unnecessary. This does not imply
that every use of four-square theory disappears from the full dependency
closure.

\subsection{Stage two: optionally change the actual derivative block}

A separate refactor could parameterize \decl{twelveSpace},
\decl{twelveTaylor}, the indicator data and their realization lemmas,
then instantiate the copy count with four. The Fourier multiplier would
have exponent $4f$ instead of $12f$, and the vacuum factor exponent
would change from twelve to four. Their rank-dependent positive height
terms decrease, but the dependent types and the corresponding product
identities still have to be checked.

In particular the original rank-one lift has block decomposition
\[
 12+4+4+(2z-20)=2z.
\]
With a four-coordinate derivative block the compatible decomposition is
\[
 4+4+4+(2z-12)=2z.
\]
\emph{The total ambient rank remains $2z$.} The complementary block
grows; it does not vanish. The source's large lower bound on $z$ leaves
room for these blocks, and their parities remain even. Rewriting all
relevant block-character and finite-cardinality identities is nevertheless
a real implementation task.\cite{rankblocks}

This stage has not been carried out. A global search-and-replace of
``12'' by ``4'' would not be justified. Nor does the local change shorten
the main lift chain $m\asymp\sqrt{x}$, improve the RH exponent, or
remove the Galois/Selmer, period-normalization or seesaw inputs.

\subsection{Files and declarations at the integration boundary}

\begin{longtable}{@{}p{.31\linewidth}p{.63\linewidth}@{}}
\toprule
Source file & Intended change or reused interface\\
\midrule\endhead
\texttt{LogTheta.lean} & Reuse the existing norm-pair integrality theorem;
no new theta isogeny theorem.\\[3pt]
\texttt{FormalParameter.lean} & Reuse the root-of-unity norm-pair lemma;
add an inert norm-minus-one selector.\\[3pt]
\texttt{LeastValuation.lean} & Generalize the cancelling product and the
unit-coefficient wrapper; expose the one-copy spanning lemma.\\[3pt]
\texttt{Couplings.lean} & Reuse the generic coefficient identity; replace
six-copy wrappers with a generic or two-copy wrapper.\\[3pt]
\texttt{Cases.lean} & Add the four-copy theorem, split out ramification,
and derive the old twelve-copy theorem by padding.\\[3pt]
\texttt{BoundAttained.lean} & No change required for stage one. Optional
parameterization for stage two.\\[3pt]
\texttt{FitTwists.lean} and\newline\texttt{TwistedValueHeights.lean}
& Stage two only: migrate the derivative and complementary blocks while
keeping total rank $2z$.\\
\bottomrule
\end{longtable}

The archive's \texttt{lean/INTEGRATION.txt} gives the actionable plan,
and \texttt{sources/source\_anchors.tsv} records exact declaration lines
and blob hashes. The candidate Lean file covers the basic product,
two-factor inverse and padding lemmas; it does \emph{not} contain the
complete selectors, case split, or downstream integration.

\section{Checks performed and what remains to verify}
\label{sec:checks}

\subsection{Source and mathematical checks}

The repository was pinned before inspection. The source inventories and
declaration anchors were compared with Git objects at that revision.
Independent reviews checked the norm-pair argument, the prime cases,
the generic coupling signature, the exposed one-copy span witness,
the unchanged error terms, and the local-to-global quantifier order.
The proposal audit included optional and future-work sections, not only
the main advertised results.

The source checker in \texttt{sources/audit\_sources.py} can repeat
the recorded Git and declaration checks against a supplied local checkout.
It performs no network access. Source matching verifies what was analyzed;
it is not a substitute for Lean elaboration or kernel replay.

\subsection{Exact finite computations}

The standard-library Python suite in \texttt{checks/} was run with exact
integer, residue-ring and rational arithmetic. It passed:
\begin{itemize}
\item 61 inert cases, including exact order $2(p+1)$, norm $-1$,
Teichm\"uller lifting, and $\zeta^p=\bar\zeta$;
\item 55 split cases, including eight cases at $2$, verifying the
two selected embeddings and the norm-minus-one congruence;
\item eight ramified cases, checking the two-norm residue obstruction,
the sum-of-two-squares solutions, and their Hensel lifts;
\item six formal exponential identities through degree ten, including
nonzero norm-weight-sum controls and the unit/inverse identity.
\end{itemize}
The test parameters are $d=7,23,31,47,71,79,103,127$.
Unramified primes range through $47$; every listed ramified prime is also
tested. Odd-prime precision is $p^7$, and split-at-two precision is
$2^{12}$. At two the script retains all root branches while lifting,
rather than incorrectly applying simple-root Hensel or choosing a branch
that may fail to extend. The JSON records witnesses and lifts.

These computations catch sign, order and small-prime mistakes. They do
not prove the infinitary local-density statements, the CM theta
integrality input, or the four-copy theorem. Those depend on the
mathematical arguments above and the existing source lemmas. The formal
series tests use arbitrary rational scalar examples to check algebra;
they do not assert that those examples are actual CM norm pairs.

\subsection{Formal status and next validation gate}

No Lean executable was available in the working environment. The included
candidate module contains no proof placeholders or added axioms, but
its syntax, typeclass inference and API uses have not been elaborated.
It should be treated as implementation scaffolding with proposed proof
terms. An uncompiled file is not formal evidence merely because it has
no visible \texttt{sorry}.

The next validation gate is to compile the local lemmas against the pinned
Lean/Mathlib versions, complete the generic wrappers and four-copy theorem,
then rebuild the unchanged twelve-copy downstream interface. Only after
that should one measure source/build effects or attempt stage two.
There is no measured threefold speedup or whole-proof size reduction in
this report; the factor of three refers to the \emph{local witness copy
count}, twelve to four.

\Needspace{13\baselineskip}
\section{Assessment}

The strongest recommendation supported by the review is a focused
arithmetic refactor: use the CM norm coordinate already present in the
integrality theorem, construct two cancelling factors at unramified primes
and three at the ramified prime, and combine the existing generic
coefficient bounds to prove four-copy sharpness.

The mathematical gain is concrete: the same local valuation budget is
achieved with fewer auxiliary copies, the universal four-squares
construction is replaced by elementary local CM choices, and an existing
one-copy result is exposed instead of discarded by padding. The safe
integration route lets this coexist with the earlier proposals while
preserving the exact determinant main term and all global normalization
interfaces. The remaining work is formalization of this replacement,
not a new conjectural analytic estimate.

\appendix
\section{Reproduction and archive contents}

The archive is self-contained for reading, building the PDF, and running
the finite checks. It does not redistribute the 803-file Lean project or
the Mathlib cache. Obtain the pinned repository separately to run source
audits or complete Lean integration.

\begin{center}\small
\begin{tabularx}{\linewidth}{@{}lX@{}}
\toprule
Path & Contents\\\midrule
\texttt{report.tex}, \texttt{report.pdf} & This report and its compiled form.\\
\texttt{README.txt} & Status, contents, and reproduction commands.\\
\texttt{lean/} & Uncompiled candidate, integration plan, attribution and license.\\
\texttt{checks/} & Exact Python suite, recorded output, detailed JSON witnesses.\\
\texttt{sources/} & Revision metadata, proposal inventory/exclusions,
declaration anchors, and offline source audit.\\
\texttt{SHA256SUMS} & SHA-256 checksums of the delivered contents.\\
\bottomrule
\end{tabularx}
\end{center}

Build the document with a standard TeX installation:
\begin{verbatim}
latexmk -pdf -interaction=nonstopmode -halt-on-error report.tex
\end{verbatim}
Run the finite checks from the archive root:
\begin{verbatim}
python3 checks/check_norm_cancellation.py --output checks/reproduced.json
\end{verbatim}
Run the source audit with the path to the ProveIt checkout:
\begin{verbatim}
python3 sources/audit_sources.py /path/to/ProveIt
\end{verbatim}
The \texttt{lean/INTEGRATION.txt} file explains the candidate-module
compilation command and the unimplemented pieces. The status markings
should be retained if the archive is forwarded.

\begin{thebibliography}{20}
\raggedright
\bibitem{public}
ProveIt, pinned revision \texttt{89d12dd4}.
\Src{Statements.lean}{\texttt{Statements.lean}} and
\Src{Main.lean}{\texttt{Main.lean}}. Public RH and Dirichlet-GRH statements
and proofs. Repository README, toolchain and manifest supply version and
attribution information.

\bibitem{verification}
Andrew V. Sutherland,
\href{https://math.mit.edu/~drew/shortgrh/}{Verification of the shorter
Lean proof of RH and GRH}, MIT, consulted 27 November 2026.
Upstream revision \texttt{72fead102f663d5c7a407964a13dc0cf9d977207}.

\bibitem{guide}
\href{https://github.com/VladimirReshetnikov/ProveIt/blob/89d12dd40424f1217fd818d84877b3da52396a3b/Analysis/GeneralizedRiemannHypothesis/docs/guide/theta_lifts_and_dirichlet_zeros.tex}{How theta lifts locate the zeros of Dirichlet $L$-functions},
the supplied human-readable guide, at the same pinned revision.

\bibitem{mainineq}
\Src{New/Reduction/MainInequality.lean}{\texttt{New/Reduction/MainInequality.lean}},
especially \decl{star_inequality} and the file's input/cancellation documentation.

\bibitem{endpoint}
\Src{New/Reduction/OneSidedBound.lean}{\texttt{New/Reduction/OneSidedBound.lean}};
\Src{New/Reduction/Zeros/Criterion.lean}{\texttt{New/Reduction/Zeros/Criterion.lean}};
\Src{New/Main.lean}{\texttt{New/Main.lean}}.

\bibitem{normpairs}
\Src{Lib/General/NumberTheory/PadicLimitsValuation.lean}{\texttt{PadicLimitsValuation.lean}},
\decl{normPairClosure} and its Teichm\"uller limit theorem;
\Src{Lib/CM/FormalGroup/FormalParameter.lean}{\texttt{FormalParameter.lean}},
\decl{mem_normPairClosure_of_pow_eq_one}.

\bibitem{logtheta}
\Src{Lib/CM/Integrality/LogTheta.lean}{\texttt{Lib/CM/Integrality/LogTheta.lean}},
especially \decl{expSubst_sub_rescale_mem}.

\bibitem{leastval}
\Src{New/Derivatives/Sharpness/LeastValuation.lean}{\texttt{LeastValuation.lean}},
especially \decl{LeastValLE}, the copy-substitution and monomial-extraction
lemmas, the six-scale constructor, and the one-copy subproof in
\decl{leastValLE_zero_of_span}.

\bibitem{couplings}
\Src{New/Derivatives/Sharpness/Couplings.lean}{\texttt{Couplings.lean}},
especially \decl{IsCouplingFamily.norm_coeff_pairExponent} and the
specialized six-copy wrappers.

\bibitem{nearpairs}
\Src{New/Derivatives/Sharpness/NearPairs.lean}{\texttt{NearPairs.lean}},
the near-pair restrictions and loss-product definitions.

\bibitem{cases}
\Src{New/Derivatives/Sharpness/Cases.lean}{\texttt{Cases.lean}},
especially \decl{gammaPrime_mul_degree_sub_le_localPrimeSumTotal_add},
the prime cases, and \decl{leastValLE_twelve}.

\bibitem{boundattained}
\Src{New/Derivatives/Sharpness/BoundAttained.lean}{\texttt{BoundAttained.lean}},
especially \decl{exists_le_norm_coeff_indicatorDatum} and
\decl{exists_rpow_neg_localPrimeSumTotal_le_norm_coeff}.

\bibitem{rankblocks}
\Src{New/Determinant/Twist/FitTwists.lean}{\texttt{FitTwists.lean}},
\decl{RankOneIndex};
\Src{New/OneLift/HeckeValues/TwistedValueHeights.lean}{\texttt{TwistedValueHeights.lean}},
rank-one test-form blocks and heights.

\bibitem{sparse}
\PropSrc{rh_sparse_theta_proof/sparse_theta_proof.tex}{\texttt{sparse\_theta\_proof.tex}}.
The other four proposal bundles are linked in Section~\ref{sec:exclusions};
the complete exclusion inventory is delivered with this report.
\end{thebibliography}
\end{document}
